\documentclass[11pt,a4paper]{article}
\usepackage[T1]{fontenc}
\usepackage[utf8]{inputenc}
\usepackage[provide=*,french]{babel}
\usepackage{lmodern,microtype}
\usepackage[a4paper,margin=23mm,headheight=15pt]{geometry}
\usepackage{amsmath,amssymb,mathrsfs,booktabs,tabularx,array,enumitem,fancyhdr,xcolor}
\usepackage[colorlinks=true,linkcolor=blue!40!black,urlcolor=blue!40!black]{hyperref}
\usepackage{xurl}
\newcommand{\MainPageCount}{63}
\newcommand{\CompanionPageCount}{96}
\newcommand{\PreservedLabelCount}{545}
\newcommand{\CurrentLabelCount}{546}
\newcommand{\NumberedResultCount}{167}
\newcommand{\CrossDocumentLinkCount}{228}
\newcommand{\InternalLinkCount}{660}

\definecolor{accent}{HTML}{173B5B}
\pagestyle{fancy}\fancyhf{}
\fancyhead[L]{\small\color{accent} P1 -- Note de révision V1.6}
\fancyhead[R]{\small 31 octobre 2026}
\fancyfoot[C]{\thepage}
\renewcommand{\headrulewidth}{0.3pt}
\setlength{\parindent}{0pt}
\setlength{\parskip}{6pt}
\setlength{\emergencystretch}{2em}
\setlist{itemsep=4pt,topsep=4pt,leftmargin=18pt}
\renewcommand{\arraystretch}{1.13}
\newcommand{\E}{\mathbb E}
\newcommand{\Law}{\mathcal L}
\newcommand{\eps}{\varepsilon}
\hypersetup{pdftitle={P1 V1.6 -- Reponse a l'audit de suivi},
pdfauthor={Brice Pouly},pdfsubject={References, mecanismes, reconstruction et relectures ciblees}}
\widowpenalty=10000\clubpenalty=10000
\begin{document}
{\Large\bfseries\color{accent} P1 V1.6 : réponse à l'audit de suivi}\\[4pt]
{\large Note de révision -- 31 octobre 2026}

\section*{1. Une consolidation ciblée de l'exposition}
La V1.6 répond au document
\nolinkurl{AUDIT_P1_V1_5_2026-10-31_FR.pdf}.
Elle conserve la scission en \textbf{deux documents associés} :
\textbf{\MainPageCount{} pages} pour l'article et
\textbf{\CompanionPageCount{} pages} pour le compagnon.
Les changements portent sur la comparaison des références, l'explication
des deux mécanismes d'erreur, les étapes de reconstruction et la navigation.
Cette édition ne revendique aucun nouvel exposant.

L'audit ne fournit pas de défaut bloquant démontré dans son périmètre.
Les trois contre-lectures supplémentaires décrites plus loin n'en ont pas
identifié non plus dans les arguments ciblés. Ces constats ne valent pas
certification exhaustive du travail. Les recommandations antérieures
déjà satisfaites ne sont pas rouvertes sans motif.

\subsection*{E1. Réunir les références et leur précision}
Une table à la fin des résultats introductifs rassemble les six
comparaisons. Elle renvoie à leurs énoncés respectifs et conserve les
conditions qui en déterminent la portée.

\small
\begin{tabularx}{\textwidth}{@{}>{\raggedright\arraybackslash}p{.33\textwidth}>{\raggedright\arraybackslash}X@{}}
\toprule
Référence ou observation & Résultat établi\\
\midrule
Poisson de moyenne \(\lambda_N\) ou \(\mu_N\)
& Équivalent explicite d'ordre \(LN^{-2/3}\).\\[5pt]
Quadratique exacte \(Q_{\mu_N,\alpha_N}\)
& Puissance optimale \(4/3\), avec marges arbitraires fixées ;
pas d'équivalent complet du reste.\\[5pt]
Quadratique nominale \(Q_{\lambda_N,\alpha_N}\)
& Majorant \(N^{-1+\eps}\) ; distance comparable au biais
\(|\mu_N-\lambda_N|\) lorsque celui-ci est non nul.\\[5pt]
Quadratique à canaux jusqu'à \(h\)
& Équivalent d'ordre \(LN^{-h/(h+1)}\) venant du premier
étage omis ; \(h\ge3\) reste fixé.\\[5pt]
Hermite à paramètres exacts
& Majorant \(N^{-4/3+\eps}\), sans minorant optimal
correspondant établi ici.\\[5pt]
Champ positionnel non signé
& Majorant \(N^{-4/5+\eps}\), avec tout poids polynomial
de degré fixé dans le compte total.\\
\bottomrule
\end{tabularx}
\normalsize

Les \textbf{\NumberedResultCount{} énoncés numérotés} de type théorème,
proposition, lemme ou corollaire conservent leur texte LaTeX exact,
leur document et leur numéro. Les \textbf{\PreservedLabelCount{} repères
de la V1.5} sont préservés. Un nouveau repère désigne la table.

\subsection*{E5. Alléger la discussion finale}
Le début de la section 19.3 renvoie désormais à cette table, puis expose
les questions restantes. Les arguments sur le test de Charlier, les
moments fixés, l'absence d'un contrôle exponentiel importable et la
géométrie des poids positionnels sont conservés. Aucun résultat ni
environnement de preuve n'est supprimé.

\clearpage
\section*{2. Deux mécanismes, deux échelles}
\subsection*{E2. Expliquer le témoin ternaire dès l'introduction}
L'introduction garde l'explication du canal \(2:3\), du marqueur de
début de run et de l'absence de famille intérieure macroscopique
\(1:2\) dans le bloc dyadique. La section 1.2 anticipe ensuite le témoin
\(10:12:15\), dont la preuve demeure en section 14.

Multiplier les fenêtres de centres \(10t,12t,15t\) par \(6,5,4\)
donne le centre commun \(60t\) et les trois pas \(6,5,4\).
Les recouvrements rationnels donnent
\[
 \rho_L=\left\lfloor\frac{L-1}{6}\right\rfloor
       +\left\lfloor\frac{L-1}{3}\right\rfloor
       +\left\lfloor\frac{L-1}{5}\right\rfloor
       -\left\lfloor\frac{L-1}{15}\right\rfloor
       =\frac{19}{30}L+O(1).
\]
En l'absence de relations supplémentaires, le poids d'un triple
compatible a l'ordre \(N^{-3+19/30}=N^{-71/30}\).
Une population d'ordre \(N\) explique donc la puissance \(41/30\).

Cette lecture géométrique ne suffit pas au minorant du cumulant.
La proposition 14.5 vérifie aussi les coordonnées privées, les signes
des prédécesseurs, la positivité du coefficient centré et les hôtes
exceptionnels. Le majorant portant sur \emph{tous} les triples utilise
les autres secteurs de la section 14 et le compagnon T10.
Le coefficient \(1/120\), sous \(N=2^L\) et \(30\mid L\), concerne
la somme des coefficients centrés d'une sous-famille de translations,
donc sa contribution à \(K_3^+\). Il n'est pas une constante du
cumulant entier, dont le raccord à \(K_3^+\) comporte le facteur six.

\subsection*{Pourquoi deux paires imposent une échelle plus grande}
Le second cumulant factoriel a l'ordre \(\alpha_N\asymp LN^{-2/3}\).
Deux interactions de paires produisent ainsi
\(\alpha_N^2\asymp L^2N^{-4/3}\), plus grand que l'échelle ternaire.
La référence exponentielle de Hermite explicite la forme du terme :
\[
 H_{\mu,\alpha}
 =\Pi_\mu*\exp_*(\alpha\Delta^{*2}/2)
 =Q_{\mu,\alpha}
  +\frac{\alpha^2}{8}\Pi_\mu*\Delta^{*4}
  +O_{\ell^1}(\alpha^3).
\]
Ce développement ne prouve pas, à lui seul, que la vraie loi possède
un reste de cette taille. La section 16 établit la non-annulation
arithmétique pertinente :
\[
 \kappa_{4,N}+3\alpha_N^2
 =24K_4^++3\alpha_N^2+O(N^{-3/2+\eps}),\qquad K_4^+\ge0.
\]
Ici \(K_4^+\) somme les coefficients centrés des quadruples connexes
compatibles. Le test de Charlier est non borné ; sa conversion en
minorant de norme utilise les moments d'ordre fixé, prouvés séparément.
Elle fournit la puissance optimale \(4/3\) avec une marge arbitraire,
sans démontrer un équivalent \(cL^2N^{-4/3}\) du reste.

La marge \(41/30-4/3=1/30\) permet le majorant quadratique.
Le minorant quartique porte sur cette référence précise. L'Hermite
contient déjà le terme de deux paires et ne reçoit pas le même minorant.

\clearpage
\section*{3. Rendre les passages de preuve explicites}
\subsection*{E3. Le bilan local de reconstruction}
L'ouverture de la section 15 affiche désormais
\[
 \begin{aligned}
 \|\Law(Z_N)-Q_{\mu_N,\alpha_N}\|_1
 &\le \|\Law(Z_N)-T_{\mathscr D_2}\|_1
      +\|T_{\mathscr D_2}-B_2\|_1\\
 &\qquad+\|B_2-Q_{\mu_N,\alpha_N}\|_1.
 \end{aligned}
\]
La classe \(\mathscr D_2\) retient les familles directes et les
familles \(e\cup U\) dont la paire dépendante \(e\) et les blocs
restants forment une somme directe après remplacement de la paire
par son espace \(S_e\). La condition porte sur la somme directe
de toute la collection, pas seulement sur des intersections deux à deux.

Le premier terme mesure la frontière de l'arrêt ; le suffixe non
développé reste une mesure conditionnelle positive. Le deuxième
contient les déficits indépendants et les déficits de paires, avec
leurs covariances signées. Une paire incompatible est autorisée dans
la classe : sa masse jointe est nulle, mais sa covariance reste dans
\(B_2\). Le troisième est la conversion du \(B_2\) complet, produit
des Bernoulli hétérogènes augmenté de \emph{toutes} les covariances.
Son coût est \(N^{-5/3+\eps}\), conformément au lemme 15.1.

Le théorème 15.13 assemble ces trois estimations. Les ordres fixes
de 3 à 256 et la queue uniforme de la proposition 15.12 restent
des entrées distinctes. Le bilan ajoute un affichage non numéroté :
les quatorze énoncés de la section et ses équations gardent leurs numéros.

\subsection*{E4. Entrer dans le compagnon, puis revenir}
Le guide du compagnon contient une table locale de cinq parcours.
Elle vise les \emph{sous-sections de preuve complète}, car certains
lemmes ou propositions portent le même numéro imprimé qu'une sous-section.

\small
\begin{tabularx}{\textwidth}{@{}>{\raggedright\arraybackslash}p{.24\textwidth}>{\raggedright\arraybackslash}p{.25\textwidth}>{\raggedright\arraybackslash}X@{}}
\toprule
Entrée de l'article & Preuve du compagnon & Retour dans l'article\\
\midrule
Proposition 9.1 & Sous-section T5.1 & Corollaire 9.2, projections courtes.\\[5pt]
Théorème 10.1 & Sous-section T6.3 & Corollaire 10.2, hôtes connexes.\\[5pt]
Théorème 13.2 & Sous-section T9.5 & Proposition 13.3, fibres entières,
puis comptage uniforme de la section 13.5.\\[5pt]
Proposition 14.3 & Sous-section T10.1 & Corollaire 14.4, paires fortement alignées.\\[5pt]
Lemme 15.2 & Sous-section T11.1 & Lemme 15.3, puis masse des primitives.\\
\bottomrule
\end{tabularx}
\normalsize

Cette navigation complète la carte de dépendances de la section 1.5.
Elle ne la remplace pas et ne modifie pas la répartition entre article
et compagnon. La branche régulière reste nécessaire à la comparaison
positionnelle ; elle n'est pas supprimée au motif que la branche
scalaire utilise maintenant toutes les marginales exactes.

\clearpage
\section*{4. Contre-lectures et non-régression}
\subsection*{S1. Trois contrôles techniques séparés}
Les contre-lectures assistées ont reconstruit les points désignés
par l'audit à partir des preuves écrites. Leurs comptes rendus
complets sont inclus dans \texttt{editorial/reviews/}.

\textbf{T5--T6 : degré et induction conditionnelle.}
Le contrôle porte sur le choix d'une paire de facteurs et d'un
diviseur prescrit, sans énumération de toutes les allocations ;
la perte explicite en degré ; l'exposition des incidences paires à
ordre de moment fixé ; les deux sens du comptage des certificats.
L'induction diminue le nombre de cibles de deux ou d'une selon le
receveur, tout en étant simultanée pour chaque nombre fixé de sources.
Le nombre de seuils est choisi après la précision et avant \(N\).

\textbf{T9 : de la base courte à la courbe commune.}
Les quatre obligations ont été suivies : extraction, auxiliaire
à deux témoins, facteur commun à toute la base, puis fibres entières
dans la grille d'offsets. La division par le pgcd polynomial précède
la spécialisation. La hauteur linéaire en \(q_0\) de l'auxiliaire est
distinguée du coût quadratique en \(q_0\) du résultant. La version
primaire de Binyamini--Cluckers--Novikov citée dans les manuscrits
confirme le caractère entier et uniforme en degré de l'entrée utilisée.

\textbf{T11 et section 15 : deux cibles et queue uniforme.}
Le contrôle distingue les deux projections injectives et les comptes
de composantes par cible. Il suit le choix de la hauteur des coefficients
avant le nombre d'équations, le retrait des petites courbes et les deux
strates de rang. Pour la queue, les quatre puissances, la marge minimale
\(1/60\), la perte \(8K_N=o(L)\), les domaines de cardinalité et la
grande frontière \(N^{-3+o(1)}\) sont reconstitués séparément.

\textbf{Conclusion de cette passe :} aucune lacune bloquante concrète
identifiée dans ces raccords ; aucune modification d'énoncé proposée.
Cela ne constitue ni une expertise humaine indépendante, ni une
formalisation, ni une certification de toutes les preuves du dossier.
Les résultats arithmétiques externes ne sont pas redémontrés ici.

\subsection*{N1. Les distinctions à conserver}
La relecture éditoriale vérifie les sept distinctions de l'audit :
moyenne exacte ou nominale ; puissance ou équivalent ; hauteur fixée
ou mobile ; paramètres fixés ou ordre croissant ; quadratique ou
Hermite ; compte scalaire, champ non signé ou marqué ; identité finie
ou estimation asymptotique. Les hashes des 167 énoncés V1.5 servent
en outre de contrôle reproductible de leur conservation textuelle.

Deux formulations sont resserrées : les poids positionnels portent
sur le \emph{compte total} ; dans l'exemple de T23, l'accord des
moments est \emph{éventuel pour chaque ordre fixé}, avec un seuil
qui peut dépendre de cet ordre. Il n'est pas simultané à tous les
ordres pour deux lois différentes à une échelle donnée.

La bibliographie scientifique reste celle de V1.5, avec 31 entrées.
Les comptes rendus privés sont des pièces éditoriales, pas des
références supplémentaires de recherche.

\clearpage
\section*{5. Vérifications et diffusion complète}
\subsection*{Contrôles reproductibles}
La commande suivante compile les deux articles, exécute les cinq
diagnostics, contrôle les références et produit cette note :
\begin{center}\texttt{python3 check\_all.py}\end{center}
Le journal structuré est \texttt{editorial/verification\_run.json}.
Le contrôle des sources vérifie la conservation des
\PreservedLabelCount{} anciens repères et des
\NumberedResultCount{} énoncés, ainsi que leur numérotation et leur texte.
Le contrôle PDF vérifie \CrossDocumentLinkCount{} renvois entre les
deux documents et \InternalLinkCount{} liens internes nommés.
Les rendus des pages modifiées sont inspectés séparément.

\small
\begin{tabularx}{\textwidth}{@{}>{\raggedright\arraybackslash}p{.30\textwidth}>{\raggedright\arraybackslash}X@{}}
\toprule
Diagnostic exécuté & Domaine et identité contrôlés\\
\midrule
Bernoulli hétérogènes
& Huit lois jointes finies, 112 contrôles : marginales, coefficients
centrés, \(B_2\), cumulants factoriels jusqu'à quatre.\\[5pt]
Arrêt héréditaire
& Toutes les classes héréditaires sur un à quatre sites contenant
l'ensemble vide, tous les états ; suffixe dépendant, déficits et
seuils finis de positivité.\\[5pt]
Coefficients de canaux
& Hauteurs jusqu'à 12, longueurs jusqu'à 96, deux extrémités ;
placements signés, normes de différences et coefficient \(1/(12e)\).\\[5pt]
Canaux et triplets
& Hauteurs 3 à 10, longueurs 5 à 150 avec \(L\ge b\) ; préfacteurs
de phase 4, 5, 6 ; triplets \(2\le u<150\), \(u<v<w<2u\) ;
nullité du témoin pour \(31\le L\le300\).\\[5pt]
Norme biquadratique
& Une instance symbolique rationnelle exacte de l'auxiliaire et de
ses annulations.\\
\bottomrule
\end{tabularx}
\normalsize

Ces diagnostics sont ceux des sources de P1, relancés pour cette
édition. Ils ne sont pas les programmes privés du nouvel audit,
qui ne sont pas joints à son PDF. Ils ne vérifient ni l'uniformité
asymptotique, ni un seuil numérique de validité des théorèmes.

\subsection*{E6. Garder la paire de manuscrits avec ses sources}
L'archive \nolinkurl{P1_V1_6_Complete_2026-10-31.zip} réunit
l'article anglais, le compagnon anglais, cette note française et
\nolinkurl{P1_V1_6_Sources_2026-10-31.zip}.
Les deux PDF scientifiques conservent les noms relatifs utilisés
par leurs liens croisés. Leurs références imprimées restent utilisables
si le lecteur PDF ne suit pas ces liens.

L'archive de sources contient les trois documents LaTeX, les commandes
de construction, la bibliographie commune, les cinq diagnostics et leurs
résultats, les tables de correspondance, les rapports ciblés et le
manifeste SHA-256. Les comptes rendus V1.4 et V1.5 sont identifiés
comme historiques. La reconstruction depuis une extraction propre
est contrôlée avant diffusion.

L'assistance de ChatGPT (GPT-6 Astra) demeure déclarée dans les deux
manuscrits. L'auteur reste responsable des résultats et de toute soumission.
\end{document}
